% !TeX root = ../main.tex
\chapter{Requirements Behind the Architecture}
\label{ch:requirements}

This chapter answers whether the architectural choices fit the original
requirements. It traces in both directions: from each requirement that shaped
the structure to the structure it produced, and from each figure in this
document back to the requirements that informed it.

\section{Requirements That Shaped the Structure}

\begin{longtable}{L{0.13\textwidth}L{0.37\textwidth}L{0.42\textwidth}}
	\caption{Requirements and the architecture they produced}\label{tab:req-to-arch}\\
	\toprule
	\textbf{ID} & \textbf{Requirement} & \textbf{Architectural consequence} \\
	\midrule
	\endfirsthead
	\toprule
	\textbf{ID} & \textbf{Requirement} & \textbf{Architectural consequence} \\
	\midrule
	\endhead

	\req{NFR-8} & A generated diagram is understandable without reading the
	code or model it came from. & Abstraction is a subsystem, not a step
	inside another one. It is the only subsystem that exists for a
	non-functional requirement. Figure~\ref{fig:pipeline}. \\
	\addlinespace
	\req{NFR-2} & A failure in the tool must not terminate GENESYS. & Two
	processes rather than one, with no shared address space.
	Figure~\ref{fig:subsystems}. \\
	\addlinespace
	\req{NFR-3} & A cancelled or failed operation leaves nothing partially
	written. & The review step is a shared filter at the end of every
	pipeline, and changes are held as objects until accepted.
	Section~\ref{sec:command}. \\
	\addlinespace
	\req{NFR-4} & The same inputs always produce the same output. & Fixed
	stage order, each stage a function of its input. This is the property the
	\pattern{Repository} alternative would have cost us,
	Section~\ref{sec:repository}. \\
	\addlinespace
	\req{NFR-6} & The projects not referencing GENESYS build without it. &
	The vendor dependency is confined to one project, so the subsystem
	boundary is also the licence boundary. Figure~\ref{fig:subsystems}. \\
	\addlinespace
	\req{CM-2} & The tool works from the schema the project actually defines.
	& The Repository Client reads the schema before the Model Mapper
	addresses anything, which fixes an ordering between two subsystems. \\
	\addlinespace
	\req{CM-5} & Entities and relationships only, no layout or formatting. &
	No diagramming subsystem exists. This follows from GENESYS being a
	\pattern{Repository} architecture, Figure~\ref{fig:repository}. \\
	\addlinespace
	\req{CA-3} & The user controls where translation starts and how far it
	goes. & Work is a bounded queue rather than an open traversal, which is
	also what keeps the \pattern{Client-Server} traffic bounded. \\
	\addlinespace
	\req{AC-1} & Generate C\# or Visual Basic. & \textsc{Strategy} inside the
	Code Generator, over one canonical form.
	Figure~\ref{fig:canonical}. \\
	\bottomrule
\end{longtable}

\section{Constraints That Removed Choices}

Several of the arrangements above were not decided by us. Recording which is
which matters, because a later reader needs to know what can be revisited.

\begin{longtable}{L{0.10\textwidth}L{0.40\textwidth}L{0.42\textwidth}}
	\caption{Constraints and what they fixed}\label{tab:constraints-arch}\\
	\toprule
	\textbf{ID} & \textbf{Constraint} & \textbf{What it fixed} \\
	\midrule
	\endfirsthead
	\toprule
	\textbf{ID} & \textbf{Constraint} & \textbf{What it fixed} \\
	\midrule
	\endhead

	\textbf{C-1} & The Client API assemblies target .NET Framework 4.8 and
	calling code must match. & The split across runtimes, and therefore the
	existence of \code{Pentateuch.Absyn} on .NET Standard 2.0. \\
	\addlinespace
	\textbf{C-3} & GENESYS loads its own Roslyn, and two versions cannot share
	a process. & That the analysis runs outside GENESYS. Constrains
	\req{FS-1}. \\
	\addlinespace
	\textbf{C-6} & Diagrams are views generated from the repository. & That no
	diagramming component is needed. \\
	\addlinespace
	\textbf{C-7} & Classes and attributes are per project and may be aliased.
	& That the schema is read at run time, fixing the order of the Repository
	Client and the Model Mapper. \\
	\addlinespace
	\textbf{C-11} & Sequencing is carried by a fixed set of behaviour
	constructs. & \textsc{Composite} for the tree and \textsc{Visitor} for
	dispatch, Section~\ref{sec:composite-visitor}. \\
	\addlinespace
	\textbf{C-12} & A state machine cannot express simultaneity. & That
	\req{AS-3} reports rather than flattens, and that the two directions of
	the activity and state pair are specified separately. \\
	\bottomrule
\end{longtable}

\section{Figures Traced to Requirements}

Each figure in this document is listed below with the requirements that
informed it, so that a reader can check a diagram against the requirement it
serves rather than taking the connection on trust.

\begin{longtable}{L{0.26\textwidth}L{0.30\textwidth}L{0.36\textwidth}}
	\caption{Figures and the requirements behind them}\label{tab:fig-trace}\\
	\toprule
	\textbf{Figure} & \textbf{Shows} & \textbf{Requirements} \\
	\midrule
	\endfirsthead
	\toprule
	\textbf{Figure} & \textbf{Shows} & \textbf{Requirements} \\
	\midrule
	\endhead

	Figure~\ref{fig:pipeline} & The filter chain and the format on each
	pipe & \req{CA-1}, \req{CA-2}, \req{NFR-8}, and risk \textbf{R-2} at the
	process boundary \\
	\addlinespace
	Figure~\ref{fig:repository} & GENESYS as a repository, with our mapper as
	one client & \req{CM-5}, \textbf{C-6} \\
	\addlinespace
	Figure~\ref{fig:layering} & Where layering holds and where it does
	not & \textbf{C-1}, \req{NFR-6}, and the fallback in \textbf{R-1} \\
	\addlinespace
	Figure~\ref{fig:subsystems} & Subsystems grouped by project and
	runtime & \req{NFR-2}, \req{NFR-6}, \textbf{C-1}, \textbf{C-3} \\
	\addlinespace
	Figure~\ref{fig:canonical} & The abstract syntax as canonical form &
	\req{AC-1}, \textbf{C-1} \\
	\bottomrule
\end{longtable}

\section{Where the Architecture Is Still Open}

The question remaining is the exchange mechanism between the two processes, risk
\textbf{R-2}. The abstract syntax fixes what is exchanged, so only the transport
is undecided, and the fallback is serialised messages over standard streams
which needs no shared state. It is the unpaid half of the cost that
\pattern{Pipe and Filter} charges.

One assumption underlies all of it. \textbf{A-1} holds that the Client API
permits creating and updating entities and not only reading them. It is
unproven, and risk \textbf{R-1} tracks it. If it fails, the Repository Client
becomes a generator of interchange files for manual import. The architecture
survives that substitution, which is the point of putting an adapter there, but
the tool's usefulness would be materially reduced.
